\begin{question}{21}{
    Voor de bevoorrading van een vooruitgeschoven eenheid (Forward Operating Base) worden
    logistieke konvooien ingezet over een risicovolle route.
    
    De tijd $T$ (in uren) voordat een konvooi wordt opgehouden door een vertraging of incident,
    is een continue kansvariabele met de volgende kansdichtheidsfunctie:
    \[
        f(t) = \begin{cases} 
                    \frac{3}{64}\cdot t^2\cdot (4-t), &\text{ als } 0 \le t \le 4, \\
                                       0, &\text{ anders.} 
                \end{cases}
    \]
}
        \subquestion{4}{
            Toon met berekeningen en een schets van de grafiek van $f$ aan dat $f$ inderdaad voldoet aan de twee voorwaarden voor een kansdichtheidsfunctie.
        }
        \answer{
            Een functie $f$ kan dienen als een kansdichtheidsfunctie als geldt:

            \textbf{Voorwaarde 1:} de functie $f$ is niet-negatief voor alle waarden van $t$, oftewel \[f(t) \ge 0.\]
            
            Op het interval $0 \le t \le 4$ geldt dat $f(t) = \frac{3}{64}\cdot t^2\cdot (4-t)$.
            In dat geval geldt dat $t^2 \ge 0$ en $4 - t \ge 0$, en we hebben te maken met een vermenigvuldiging van niet-negatieve termen, dus $f(t) \ge 0$.
            Buiten het interval $0 \le t \le 4$ is $f(t) = 0$, dus is ook voldaan aan $f(t) \ge 0$. \rubric{2}

            \textbf{Voorwaarde 2:} de oppervlakte onder de grafiek van $f$ is gelijk aan 1, oftewel \[\int_{-\infty}^{\infty} f(t) \dt = 1.\]         
            
            Voor de tweede voorwaarde moeten we dus checken of de oppervlakte onder de grafiek van $f$ daadwerkelijk gelijk is aan 1:
            \begin{align*}
                \int_{-\infty}^{\infty} f(t)\dt   &= \int_{0}^{4} \frac{3}{64}\cdot t^2 \cdot (4 - t)\dt \\
                                                    &= 1. \rubric{2}
            \end{align*} 
            Hiermee is aangetoond dat de gegeven functie $f$ voldoet aan de twee voorwaarden voor kansdichtheidsfuncties.
        }

        \subquestion{8}{
            Bereken de verwachtingswaarde $E(T)$ en de standaardafwijking $\sigma(T)$ met behulp van de algemene formules op het formuleblad.
        }
        \answer{
            De verwachtingswaarde $E(T)$ berekenen we door middel van
            \begin{align*}
                E(T) = \int_{-\infty}^{\infty} t\cdot f(t)\dt &= \int_{0}^{4} t \cdot \frac{3}{64}\cdot t^2 \cdot (4 - t)\dt \\
                                                                &= \int_{0}^{4} \frac{3}{64}\cdot t^3 \cdot (4 - t)\dt \\   
                                                                &= 2.4. \rubric{3}          
            \end{align*}
            De standaardafwijking $\sigma(T)$ berekenen we door eerst de variantie $\Var{T}$ te bepalen, en van het resultaat de wortel te nemen:
            \begin{align*}
                \Var{T} = \int_{-\infty}^{\infty} \left(t-E(T)\right)^2 \cdot f(t)\dt    &= \int_{0}^{4} (t-2.4)^2 \cdot \frac{3}{64}\cdot t^2 \cdot (4 - t)\dt \\
                                                                                &= 0.64. \rubric{3}
            \end{align*} 
            \begin{align*}
                \sigma(T) = \sqrt{\Var{T}} = \sqrt{0.64} = 0.8. \rubric{2}
            \end{align*}
        }

        \subquestion{4}{
            Bereken de kans dat een willekeurig konvooi binnen tweeënhalf uur een vertraging of incident te verduren krijgt.
        }
        \answer{
            Hiervoor moeten we de kans $P(T < 2.5)$ berekenen:
            \[
                P(T < 2.5) = \int_{-\infty}^{2.5} f(t) \dt = \int_{0}^{2.5} \frac{3}{64} \cdot t^2 \cdot (4-t) \dt \approx 0.5188. \rubric{4}
            \]
        }

        \subquestion{5}{
            Stel dat op een gegeven moment vier logistieke konvooien onafhankelijk van elkaar vertrekken richting de FOB.
            De bevoorradingsoperatie is succesvol als er minstens twee konvooien langer dan 1.5 uur onverstoord blijft rijden. 
            Hoe groot is de kans op een succesvolle bevoorradingsoperatie?
        }
        \answer{
            Laat $Y$ het aantal konvooien zijn dat langer dan 1.5 uur onverstoord blijft rijden.
            $Y$ is binomiaal verdeeld met parameters $n = 4$ en $p = 1 - 0.5188 = 0.4812$ (de kans dat een konvooi langer dan 2.5 uur onverstoord blijft).\rubric{2}
            De kans op succes is de kans dat minstens twee konvooien deze tijdsduur haalt:
            \begin{align*}
                P(Y \ge 2)  &= 1 - P(Y \leq 1) \\
                            &= 1 - \binomcdf(n=4, p=0.4812, k=1) \\
                            &\approx 0.6588. \rubric{3}
            \end{align*}
            De kans op een succesvolle operatie is hiermee ongeveer \SI{65.88}{\percent}.
        }
\end{question}